%==============================================================================
%  ANEXO OBLIGATORIO: relación del trabajo con los ODS
%  Normativa de TFG y TFM de la UPV (BOUPV 17/09/2025), art. 10.3 y anexo I.
%  La ETSINF pide que vaya al final de la memoria.
%
%  Marca con \ods la casilla de cada fila (una por fila) y escribe debajo la
%  descripción de los ODS con relación más alta.
%==============================================================================
\chapter{Relación del trabajo con los Objetivos de Desarrollo Sostenible de la Agenda 2030}
\label{anexo:ods}

\newcommand{\ods}{$\times$}

Grado de relación del trabajo con los Objetivos de Desarrollo Sostenible (ODS).

\begin{table}[H]
  \centering
  \small
  \begin{tabular}{@{}l c c c c@{}}
    \toprule
    \textbf{Objetivos de Desarrollo Sostenible} & \textbf{Alto} & \textbf{Medio} & \textbf{Bajo} & \textbf{No procede} \\
    \midrule
    ODS 1. Fin de la pobreza.                         & & & & \\
    ODS 2. Hambre cero.                               & & & & \\
    ODS 3. Salud y bienestar.                         & & & & \\
    ODS 4. Educación de calidad.                      & & & & \\
    ODS 5. Igualdad de género.                        & & & & \\
    ODS 6. Agua limpia y saneamiento.                 & & & & \\
    ODS 7. Energía asequible y no contaminante.       & & & & \\
    ODS 8. Trabajo decente y crecimiento económico.   & & & & \\
    ODS 9. Industria, innovación e infraestructuras.  & & & & \\
    ODS 10. Reducción de las desigualdades.           & & & & \\
    ODS 11. Ciudades y comunidades sostenibles.       & & & & \\
    ODS 12. Producción y consumo responsables.        & & & & \\
    ODS 13. Acción por el clima.                      & & & & \\
    ODS 14. Vida submarina.                           & & & & \\
    ODS 15. Vida de ecosistemas terrestres.           & & & & \\
    ODS 16. Paz, justicia e instituciones sólidas.    & & & & \\
    ODS 17. Alianzas para lograr objetivos.           & & & & \\
    \bottomrule
  \end{tabular}
\end{table}

\section*{Descripción de la alineación del TFG con los ODS con un grado de relación más alto}

% Aquí va tu texto (tantas páginas como haga falta). Lo decides tú; como
% punto de partida, los que más encajan con una BCI de habla suelen ser el
% ODS 3 (salud y bienestar) y el ODS 10 (reducción de las desigualdades:
% accesibilidad para personas con discapacidad), y algo el ODS 9 (innovación).
\textit{[Pendiente de redactar.]}
